\uuid{eoM7}
\titre{Éléments simples}
\auteur{bodin ; hanani ; mzari}
\organisation{exo7}
\licence{CC BY-NC-SA 4.0}
\datecreate{2020-01-20}
\niveau{L1}
\module{Algèbre}
\chapitre{Polynôme, fraction rationnelle}
\sousChapitre{Fraction rationnelle}
\difficulte{1}
\theme{Fractions rationnelles, éléments simples}
\exo7qcmid{lille-1.polynomes.025}

\contenu{
\begin{question}
\qclassification{105.04}
\qtags{motcle=éléments simples}
Quelles sont les affirmations vraies ?
\begin{answers}
    \bad{Les éléments simples sur $\Cc$ sont de la forme $\frac{a}{X-\alpha}$, $a,\alpha \in \Cc$.}

    \bad{Les éléments simples sur $\Cc$ sont de la forme $\frac{a}{(X-\alpha)^k}$, $a,\alpha \in \Rr$, $k\in\Nn^*$.}

    \good{Les éléments simples sur $\Rr$ peuvent être de la forme $\frac{a}{(X-\alpha)^k}$, $a,\alpha \in \Rr$.}

    \bad{Les éléments simples sur $\Rr$ peuvent être de la forme $\frac{aX+b}{X-\alpha}$, $a,b,\alpha \in \Rr$.}

\end{answers}
\begin{explanations}
Sur $\Cc$ les éléments simples sont de la forme $\frac{a}{(X-\alpha)^k}$, $a,\alpha \in \Cc$, $k\in\Nn^*$.
Sur $\Rr$ les éléments simples sont de la forme $\frac{a}{(X-\alpha)^k}$, $a,\alpha \in \Rr$, $k \in \Nn^*$ ou bien
$\frac{aX+b}{(X^2+\alpha X+\beta)^k}$, $a,b,\alpha,\beta \in \Rr$, $k \in \Nn^*$ avec
$X^2+\alpha X+\beta$ sans racines réelles.
\end{explanations}
\end{question}
}
